\documentclass[12pt]{article}
\usepackage[utf8]{inputenc}
\usepackage[T1]{fontenc}
\usepackage{hyperref}
\title{Tondano Language Resources for K-12 Curriculum Development}
\author{Compiled by Claude Sonnet 5 (claude-sonnet-5) for Andr\'as Kornai}
\date{}
\begin{document}
\maketitle
\section{Introduction}
Tondano (Toundano; ISO 639-3 tdn) is a Minahasan (Austronesian) language of North Sulawesi, Indonesia, spoken mainly by older adults around the town of Tondano in Minahasa regency. Search results gave a 2010 census count of about 9,560 speakers, which differs from the roughly 92,000 in the survey description. Omniglot states that an orthography and school materials have been developed since 2001; the search did not show those materials themselves.
\section{Primary Educational Resources}
No specific primers or textbooks were found in search results. Omniglot's mention of school materials since 2001 (see below) is unverified as to titles.
\section{Linguistic and Reference Materials}
\begin{itemize}
\item Timothy C. Brickell, \emph{Tondano (Toundano): A Grammar Sketch of an Endangered Minahasan Language}, Routledge 2022: \url{https://routledge.com/Tondano-Toundano-A-Grammar-Sketch-of-an-Endangered-Minahasan-Language/Brickell/p/book/9781032253022}
\item Taylor \& Francis listing of the same book: \url{https://taylorfrancis.com/books/9781138549975}
\item WorldCat record: \url{https://search.worldcat.org/zh-cn/title/1352972219}
\item University of Bamberg research record (Toulour dialect): \url{https://fis.uni-bamberg.de/handle/uniba/97614}
\item University of Bamberg research data entry: \url{https://fis.uni-bamberg.de/entities/researchdata/29644fd2-f225-4a54-af08-525fe011313d}
\item Omniglot page on Tondano (writing and publications): \url{https://omniglot.com/writing/tondano.htm}
\end{itemize}
\section{Religious and Community Resources}
None found.
\section{Digital and Audio Resources}
No standalone audio or video resources were found. The Bamberg research data entry above may hold language data.
\section{Curriculum Development Context}
Brickell's grammar and the Bamberg data are the main foundations. Omniglot reports an orthography and school materials since 2001, which should be followed up locally. No dictionary was confirmed.
\end{document}